(1) 波速 \(c=\sqrt{T/\mu}\)、波数 \(k=\omega/c\)。将弦端取为 \(x=0\)，入射与反射场为 \(y=A[e^{i(kx-\omega t)}+r e^{i(-kx-\omega t)}]\)。端部力平衡 \(T y_x(0,t)+b y_t(0,t)=0\) 给 \(Tk(1-r)-b\omega(1+r)=0\)，因此 \(r=(Z-b)/(Z+b)\)，其中 \(Z=T/c=\sqrt{T\mu}\)。
(2) 对相同弦中的相向行波，平均功率之比为反射振幅平方 \(R=|r|^2=(Z-b)^2/(Z+b)^2\)。端部耗散比例 \(A_{\rm abs}=1-R=4Zb/(Z+b)^2\ge0\)，且 \(R+A_{\rm abs}=1\)；当 \(b=Z\) 时 \(r=0\)，负载机械阻抗匹配，理想条件下入射功率被完全吸收。
(3) 对有限固定端弦写 \(y(x,t)=\sin(kx)e^{-i\omega t}\)，\(\omega=ck\)。在 \(x=L\) 施加同样的条件 \(T y_x(L)+b y_t(L)=0\) 得 \(Tk\cos(kL)-i\omega b\sin(kL)=0\)，即 \(\cot(kL)=i b/Z\)（对非零 \(k\) 且 \(\sin(kL)\ne0\)）。允许复频率由此确定，纯实频率通常不能满足有耗端条件。
(4) \(b\to0\) 时 \(r\to+1\)，自由端位移同相反射，有限弦边界为 \(y_x(L)=0\)；\(b\to\infty\) 时 \(r\to-1\)，刚性固定端反相反射。\(b=Z\) 时理想有限弦无从端部返回的波，离散共振峰不再形成。采用 \(e^{-i\omega t}\) 时耗散模必须有 \(\Im\omega<0\)，才能随时间衰减。
